% =====================================================================
% 系统开发工具基础 · 实验报告（第 5--8 题）
%
% 编译方式（在 WSL 中）：
%   latexmk -xelatex -halt-on-error 实验报告.tex
%
% 依赖安装（首次需要）：
%   sudo apt install -y texlive-latex-base texlive-latex-recommended \
%                        texlive-latex-extra latexmk texlive-lang-chinese \
%                        texlive-xetex fonts-noto-cjk poppler-utils
% =====================================================================
\documentclass[UTF8,fontset=ubuntu]{ctexart}

\usepackage{amsmath}
\usepackage{booktabs}
\usepackage{geometry}
\usepackage{listings}
\usepackage[hidelinks]{hyperref}

\geometry{a4paper, margin=2.5cm}

\lstset{
  basicstyle=\ttfamily\small,
  breaklines=true,
  frame=single,
  framerule=0.4pt,
  tabsize=4,
  showstringspaces=false,
  columns=fullflexible,
  keepspaces=true,
}

\title{\textbf{系统开发工具基础 · 实验报告（第 5--8 题）}}
\author{姓名：李云龙 \quad 学号：25030021043}
\date{\today}

\begin{document}

\maketitle

\begin{center}
\begin{tabular}{@{}p{3cm}p{11.5cm}@{}}
\toprule
姓名 & 李云龙 \\
学号 & 25030021043 \\
实验日期 & \today \\
实验环境 & Windows 11 + WSL2（Ubuntu 26.04 LTS）；bash；git；Python 3.14 + pytest 9.1 + ruff 0.16 \\
\bottomrule
\end{tabular}
\end{center}

\begin{quote}
\textbf{说明：}本报告中各题“运行结果”均为在 WSL 中实际运行的输出；第 5 题的 Ctrl-Z / bg / jobs 等为终端交互操作，其产物（日志文件）与关键输出一并列出。
\end{quote}

\section{实验目的}
\begin{enumerate}
\item 掌握进程信号（SIGTERM/SIGINT）、trap 清理机制与 Shell 任务控制（挂起/后台/取 PID）。
\item 掌握编辑器语言服务器的跳转定义、查找引用与语义重命名，以及 ruff/pytest 的开发反馈闭环。
\item 掌握用调试器（pdb/IDE）定位逻辑缺陷的方法与最小修复原则。
\item 掌握“先测量再优化”的性能分析方法（time / cProfile）与公平对比（同输入、多次取中位数）。
\end{enumerate}

\section{实验五：控制一个可清理的后台任务（q05）}

\subsection{完整脚本（q05/worker.sh）}
\begin{lstlisting}[language=bash]
#!/usr/bin/env bash
trap 'echo CLEAN_EXIT >> cleanup.log; exit 0' TERM INT
n=0
while true; do
  echo "$n"
  n=$((n+1))
  sleep 1
done
\end{lstlisting}

赋予执行权限，前台启动并把 stdout/stderr 分别写入日志：
\begin{lstlisting}[language=bash]
cd q05
chmod +x worker.sh
./worker.sh > stdout.log 2> stderr.log
\end{lstlisting}

\subsection{交互控制步骤}
\begin{enumerate}
\item \textbf{Ctrl-Z} 挂起任务（前台进程转入停止状态，Shell 提示 \lstinline|[1]+ Stopped|）；
\item \lstinline|bg| 让任务在后台继续运行；
\item \lstinline|jobs| 确认状态：\lstinline|[1]+ Running|；
\item \lstinline|pid=$(jobs -p)| 取得任务 PID（由 Shell 管理，不手工抄写）；
\item \lstinline|kill -TERM "$pid"| 发送 SIGTERM；\lstinline|wait| 等待任务结束；
\item \lstinline|cat cleanup.log| 确认出现 CLEAN\_EXIT。
\end{enumerate}

\subsection{运行结果}
\textbf{cleanup.log}：
\begin{lstlisting}
CLEAN_EXIT
\end{lstlisting}

\textbf{stdout.log}（73 行，数字 0--72，头尾示例）：
\begin{lstlisting}
0
1
...
71
72
\end{lstlisting}

\textbf{stderr.log}：为空（程序没有错误输出）。

\subsection{要点分析}
\begin{itemize}
\item \lstinline|trap '...' TERM INT| 为进程注册信号处理：收到 SIGTERM/SIGINT 时先执行清理（追加 CLEAN\_EXIT 日志）再 \lstinline|exit 0|，实现\textbf{优雅退出}；\lstinline|kill -9| 由内核直接强杀，信号不可捕获，清理代码永远不会执行，故题目禁止。
\item Ctrl-Z / bg / jobs / jobs -p 构成 Shell \textbf{任务控制}标准链路：挂起 -> 后台 -> 查状态 -> 取 PID。用 \lstinline|jobs -p| 而非手抄 PID，避免抄错或拿到其它进程的 PID。
\item stdout 与 stderr 分离重定向便于日志分类排查：正常输出进 stdout.log，错误信息进 stderr.log。
\end{itemize}

\section{实验六：语义重构与本地开发反馈（q06）}

\subsection{文件与最终代码}
\textbf{math\_utils.py}：
\begin{lstlisting}[language=Python]
def calculate_total(price: float, count: int) -> float:
    return price * count
\end{lstlisting}

\textbf{app.py}：
\begin{lstlisting}[language=Python]
from math_utils import calculate_total

print(calculate_total(12.5, 4))
\end{lstlisting}

\textbf{test\_math\_utils.py}：
\begin{lstlisting}[language=Python]
from math_utils import calculate_total


def test_calculate_total():
    assert calculate_total(12.5, 4) == 50.0
\end{lstlisting}

\subsection{语言服务器演示步骤（VS Code + Pylance / VS + Python 扩展）}
\begin{enumerate}
\item \textbf{跳转到定义}：在 app.py 的 \lstinline|calculate_total(...)| 上按 \lstinline|F12|，跳到 math\_utils.py 中的定义；
\item \textbf{查找引用}：\lstinline|Shift+F12|，列出定义与全部使用位置（app.py 的调用、测试中的调用）；
\item \textbf{重命名符号}：\lstinline|F2| 输入 calculate\_total，语言服务器按符号身份同步修改定义与所有引用（本项目 3 处）；
\item \textbf{ruff 反馈}：临时在 app.py 加入 \lstinline|import os|（未使用）-> \lstinline|ruff check .| 报 \lstinline|F401 unused import| -> 删除该行 -> 重新检查通过。
\end{enumerate}

\subsection{运行结果}
\begin{lstlisting}
$ ruff check .
All checks passed!

$ python3 -m pytest -q
.                                                                        [100%]
1 passed in 0.07s
\end{lstlisting}

\subsection{要点分析}
\begin{itemize}
\item 语言服务器的重命名是\textbf{语义级}操作（按符号身份定位），不是文本替换，不会漏改引用、也不会误改注释或字符串中恰好相同的名字；
\item ruff 是静态检查工具，无需运行程序即可发现未使用导入（F401）、未定义变量等“代码异味”，与语言服务器、pytest 一起构成“编辑 -> 检查 -> 测试”的本地开发反馈闭环。
\end{itemize}

\section{实验七：用调试器定位归并排序缺陷（q07）}

\subsection{缺陷代码（merge\_sort\_buggy.py，题目给定）}
\begin{lstlisting}[language=Python]
def merge(left, right):
    result = []
    i = j = 0
    while i < len(left) and j < len(right):
        if left[i] <= right[j]:
            result.append(left[i]); i += 1
        else:
            result.append(right[i]); j += 1  # 此处存在缺陷
    return result + left[i:] + right[j:]


def merge_sort(a):
    if len(a) <= 1:
        return a
    m = len(a) // 2
    return merge(merge_sort(a[:m]), merge_sort(a[m:]))


print(merge_sort([3, 1, 4, 1, 5, 9, 2, 6]))
\end{lstlisting}

\subsection{运行复现（结果错误）}
\begin{lstlisting}
$ python3 merge_sort_buggy.py
[1, 5, 4, 3, 4, 5, 6, 9]        # 正确应为 [1, 1, 2, 3, 4, 5, 6, 9]
\end{lstlisting}

\subsection{调试定位（pdb）}
\begin{lstlisting}
$ python3 -m pdb merge_sort_buggy.py
(Pdb) break 8                     # 在 else 分支处设断点
(Pdb) continue
(Pdb) p i, j, left[i], right[j]   # 观察指针与两侧元素
\end{lstlisting}

单步观察发现：当 $right[j] < left[i]$ 进入 else 分支时，代码用左指针 $i$ 索引右列表 $right[i]$，取值错乱（$i$ 与 $j$ 混用）。

\subsection{最小修复与验证}
修复：else 分支 \lstinline|right[i]| -> \lstinline|right[j]|（仅改一处，不重写算法、不用 sorted）。
\begin{lstlisting}
$ python3 merge_sort.py
[1, 1, 2, 3, 4, 5, 6, 9]

$ python3 -m pytest -q            # test_merge_sort.py：给定输入 + 含重复元素
..                                                                       [100%]
2 passed in 0.07s
\end{lstlisting}

\subsection{要点分析}
\begin{itemize}
\item 先运行确认“结果错误”，再用调试器定点观察 $i$、$j$、$left[i]$、$right[j]$ 四个变量，指针混用立刻现形——比到处加 print 更高效；
\item \textbf{最小修复原则}：只改出错的一处下标，保持算法结构不变，并用两个测试（含重复元素）做回归保护。
\end{itemize}

\section{实验八：先测量，再优化慢速词频程序（q08）}

\subsection{文件}
\textbf{generate\_words.py}（可重复生成词表，固定随机种子）：
\begin{lstlisting}[language=Python]
import random

random.seed(20260831)
vocab = [f"w{i}" for i in range(1000)]
with open("words.txt", "w", encoding="utf-8") as f:
    f.write(" ".join(random.choice(vocab) for _ in range(30000)))
\end{lstlisting}

\textbf{wordfreq.py}（原版，去重用 list + in）：
\begin{lstlisting}[language=Python]
words = open("words.txt", encoding="utf-8").read().split()
unique = []
for word in words:
    if word not in unique:
        unique.append(word)
print("count=", len(unique))
\end{lstlisting}

\textbf{wordfreq\_fast.py}（优化版，set 去重）：
\begin{lstlisting}[language=Python]
words = open("words.txt", encoding="utf-8").read().split()
unique = set(words)
print("count=", len(unique))
\end{lstlisting}

\subsection{测量与分析方法}
\begin{lstlisting}[language=bash]
cd q08
python3 generate_words.py                       # 生成 words.txt（30000 词）
time python3 wordfreq.py                        # 原版，跑 2 次
python3 -m cProfile -s cumulative wordfreq.py   # 热点分析
time python3 wordfreq_fast.py                   # 优化版，跑 2 次
\end{lstlisting}

cProfile 结果：热点集中在 \lstinline|wordfreq.py:1| 的去重循环（list 的 \lstinline|in| 成员判断，tottime 0.074s），而非文件读取/分词。

\subsection{实测数据（同一 words.txt，各跑 2 次）}
\begin{center}
\begin{tabular}{@{}lccc@{}}
\toprule
版本 & run1 & run2 & 中位数 \\
\midrule
原版（list + in） & 0.099s & 0.093s & 0.096s \\
优化版（set 去重） & 0.026s & 0.019s & 0.023s \\
\bottomrule
\end{tabular}
\end{center}

\textbf{加速比 $\approx 0.096 / 0.023 \approx 4.3$ 倍}；最终 count 均为 1000，保持不变（未写死词表大小）。

\subsection{要点分析}
\begin{itemize}
\item \textbf{先测量再优化}：cProfile 证明热点是去重循环的成员判断，而不是文件读取，避免优化错方向；
\item 复杂度：list 的 \lstinline|in| 是线性扫描（整体 $O(n^2)$），set 基于哈希表（整体 $O(n)$）；
\item 对比方法：\textbf{相同输入}、每版本\textbf{跑 2 次取中位数}再算加速比，抵消系统噪声，结论更可信。
\end{itemize}

\section{遇到的问题与解决}
\begin{center}
\begin{tabular}{@{}p{5cm}p{5.2cm}p{5cm}@{}}
\toprule
问题 & 原因 & 解决 \\
\midrule
第 7 题归并结果错乱 & merge 的 else 分支 right[i] 下标混用（应为 right[j]） & pdb 断点观察指针后一行修复，测试回归 \\
第 8 题原版去重慢 & list 的 in 为线性扫描，整体 $O(n^2)$ & 改用 set 去重（$O(n)$），count 不变 \\
第 6 题未使用导入检查 & import os 未被使用 & 用 ruff 静态检查发现 F401 并删除 \\
环境缺少 ruff & Ubuntu 26.04 的 apt 无 ruff 包 & pip3 install --break-system-packages ruff pytest \\
\bottomrule
\end{tabular}
\end{center}

\section{实验总结}
通过本次实验，我掌握了进程信号与 trap 优雅退出机制，以及 Ctrl-Z / bg / jobs / jobs -p 的任务控制链路；体验了编辑器语言服务器的跳转定义、查找引用与语义重命名，并用 ruff + pytest 建立了本地开发反馈闭环；学会用调试器定点观察变量来定位“指针/下标混用”类缺陷并做最小修复；理解了“先测量再优化”的分析方法，通过 cProfile 定位热点后把去重从 $O(n^2)$ 优化到 $O(n)$，实测加速约 4.3 倍，并学会了用中位数做公平的性能对比。

\section*{附录：文件清单}
\begin{lstlisting}
w2/README.md                      第 5--8 题说明（步骤、原理、实测数据）
w2/q05/worker.sh                  第 5 题脚本（trap TERM/INT 清理）
w2/q05/cleanup.log / stdout.log / stderr.log   第 5 题运行产物
w2/q06/math_utils.py / app.py / test_math_utils.py   第 6 题（已重命名 calculate_total）
w2/q07/merge_sort_buggy.py        第 7 题缺陷版
w2/q07/merge_sort.py              第 7 题修复版
w2/q07/test_merge_sort.py         第 7 题测试（2 个用例）
w2/q08/generate_words.py / wordfreq.py / wordfreq_fast.py / words.txt   第 8 题
实验报告.md / 实验报告.tex / 实验报告.pdf    本报告
\end{lstlisting}

\end{document}
